% cau-hoi.tex: gói 1 của bộ GAME Jelly Trail Academy (dulieu=jelly), 20 câu. SQL nằm trong solution.sql, nối với câu bằng mã câu.
% Biên dịch trên Overleaf: Menu > Compiler > XeLaTeX. Hướng dẫn: HUONG_DAN.md.
% Phạm vi: Module 1 Retrieval (9 câu), Module 2 NULL & Expressions (9 câu), 2 case study cấp Cơ bản.
% Thuật ngữ, giọng văn: style/STYLE-jelly.md. Bản đồ câu: goi-cau-hoi/jelly/BAN-DO-100-CAU.md (gói 1).
\documentclass[11pt,a4paper]{article}
\usepackage{aio-sql-cauhoi}
% Hai dòng vá XeLaTeX (cong-cu/va-xelatex.tex), không sửa .sty:
% 1) font chữ đơn cách có đủ chữ Việt; 2) \erdkhoa trong \begin{erd} vẽ thành node để không mất chữ.
\setmonofont{NotoSansMono}[Extension=.ttf, UprightFont=*-Regular, BoldFont=*-Bold, Scale=0.86]
\renewcommand{\erdkhoa}[1]{\node[font=\small, color=cqMuted, anchor=north west, align=left, text width=14cm] at ([yshift=-1mm]current bounding box.south west) {\textbf{Khoá:} #1};}

% ================================================================ KHAI BÁO DÙNG CHUNG
% Câu hỏi chỉ được dùng các mã khai báo ở đây.
% ---------------------------------------------------------------- bộ dữ liệu
\khaibaodulieu{jelly}{GAME · Jelly Trail Academy}
  {Game giải đố ghép kẹo trên di động (iOS, Android, 8 thị trường): người chơi, phiên chơi, lượt chơi màn, mua trong ứng dụng, quảng cáo, chi phí thu hút; 60 ngày từ 2026-07-01 tới 2026-08-29.}

% ---------------------------------------------------------------- domain (6 khoá chuẩn, ánh xạ theo STYLE-jelly.md mục G)
\khaibaodomain{customer}{Người chơi}{Hoạt động, giữ chân, phễu hướng dẫn, nhóm cài đặt.}
\khaibaodomain{sales}{Mua trong ứng dụng}{Giao dịch IAP, gói vật phẩm, người chơi trả tiền.}
\khaibaodomain{marketing}{Thu hút và sự kiện}{Kênh, chiến dịch, chi phí, sự kiện LiveOps, thử nghiệm A/B.}
\khaibaodomain{operations}{Vận hành game}{Màn chơi, độ khó, booster, phiên chơi, phiên bản.}
\khaibaodomain{finance}{Doanh thu}{Tổng doanh thu, ARPDAU, ARPPU, CPI, ROAS.}
\khaibaodomain{after_sales}{Chất lượng sau phát hành}{Sự cố, hoàn tiền, đối soát sau bản vá.}

% ---------------------------------------------------------------- khái niệm (từ STYLE-jelly.md mục E, chỉ các khái niệm gói này dùng)
% \khaibaokhainiem{ma}{Tên}{Định nghĩa}{Công thức}{Trong bộ dữ liệu}

% --- Người chơi, phiên chơi, phiên bản
\khaibaokhainiem{dau}{DAU (số người chơi hoạt động trong ngày)}
  {DAU là số người chơi khác nhau có mở game ít nhất một lần trong một ngày.}
  {số user\_id khác nhau có hoạt động trong ngày}
  {bảng \code{player_daily}: mỗi dòng là một người chơi trong một ngày có hoạt động; \code{daily_kpi.dau} là số tổng hợp sẵn, khớp với cách đếm này.}
\khaibaokhainiem{phien-choi}{Phiên chơi và thời lượng phiên}
  {Phiên chơi là một lần người chơi mở game cho tới khi rời đi. Thời lượng phiên là số phút của phiên đó.}
  {thời lượng phiên bình quân = tổng duration\_minutes / số phiên}
  {\code{sessions}; mã phiên có hai dạng, phiên hướng dẫn có \code{entry_screen} = \code{'tutorial'}.}
\khaibaokhainiem{phien-ban-luc-choi}{Phiên bản lúc chơi và lúc cài}
  {Phiên bản lúc chơi là phiên bản ứng dụng người chơi dùng trong ngày hoạt động; người chơi có thể cập nhật sau khi cài.}
  {}
  {\code{player_daily.app_version}, \code{sessions.app_version}, \code{level_attempts.app_version} là phiên bản lúc chơi; \code{users.app_version_install} là phiên bản lúc cài. Bản 2.1.0 có từ 30/7, 2.1.1 từ 15/8.}
\khaibaokhainiem{ty-le-su-co}{Tỷ lệ sự cố trên 1.000 phiên}
  {Tỷ lệ sự cố đo số lần ứng dụng tự thoát do lỗi trên mỗi 1.000 phiên chơi.}
  {tổng crash\_count / tổng sessions × 1000}
  {\code{player_daily.crash_count}, \code{player_daily.sessions}. Số lượt chơi có \code{result} = \code{'crash'} là nguồn đếm khác, không khớp với \code{crash_count}.}

% --- Doanh thu
\khaibaokhainiem{doanh-thu-iap}{Doanh thu IAP}
  {Doanh thu IAP là tiền người chơi trả để mua gói vật phẩm ảo trong ứng dụng.}
  {tổng price\_usd của giao dịch không hoàn tiền}
  {bảng \code{purchases}: mỗi dòng là một giao dịch, giá theo USD; trong kỳ chưa có giao dịch hoàn tiền (\code{is_refund} = 0).}
\khaibaokhainiem{doanh-thu-quang-cao}{Doanh thu quảng cáo}
  {Doanh thu quảng cáo là tiền mạng quảng cáo trả cho game khi người chơi xem quảng cáo.}
  {tổng revenue\_usd}
  {\code{ad_events} (từng lần xem, 5 chữ số thập phân); \code{player_daily.ad_revenue} là tổng theo ngày đã làm tròn 4 chữ số.}
\khaibaokhainiem{nguoi-choi-tra-tien}{Người chơi trả tiền}
  {Người chơi trả tiền là người có ít nhất một giao dịch mua trong ứng dụng trong kỳ xét.}
  {số user\_id khác nhau có giao dịch}
  {\code{purchases}; một người có thể có nhiều giao dịch. Theo ngày, người chơi trả tiền là người có \code{player_daily.iap_revenue} lớn hơn 0.}
\khaibaokhainiem{arpdau}{ARPDAU}
  {ARPDAU là doanh thu bình quân trên một người chơi hoạt động trong ngày.}
  {(doanh thu IAP + doanh thu quảng cáo) / DAU}
  {\code{daily_kpi.arpdau} làm tròn 4 chữ số. ARPDAU của nhiều ngày tính bằng tổng doanh thu chia tổng DAU, không lấy trung bình cột \code{arpdau}.}

% --- Thử nghiệm
\khaibaokhainiem{thu-nghiem-ab}{Thử nghiệm A/B theo nhánh được chia}
  {Thử nghiệm A/B so sánh hai nhánh người chơi được chia ngẫu nhiên, mỗi nhánh nhận một cách thiết kế khác nhau.}
  {chỉ số của nhánh thử − chỉ số của nhánh đối chứng}
  {\code{ab_assignments}: thử nghiệm \code{extra_moves_offer_timing}, nhánh đối chứng \code{control_after_3_fails}; chỉ số chính tính trên mọi người được chia, kể cả người chưa tiếp xúc (\code{exposed_at} rỗng).}

\begin{document}
{\LARGE\bfseries Gói câu hỏi SQL · GAME Jelly Trail Academy · gói 1\par}
\medskip
{\color{cqMuted}Bộ dữ liệu \texttt{jelly} · Module 1 Retrieval (9 câu) · Module 2 NULL \& Expressions (9 câu) · 2 case study Cơ bản · chuẩn \texttt{aio-sql-cauhoi} v0.4\par}
\bigskip
Gói chờ giảng viên duyệt. Khung xanh nét đứt (\emph{Thiết kế ERD}) là ghi chú cho người duyệt. Các khung Truy vấn khởi đầu, Lời giải và Bài sai lấy từ \code{solution.sql}; học viên chỉ thấy truy vấn khởi đầu.

\tableofcontents


% ################################################################ MODULE 1 · RETRIEVAL

% ================================================================ R1 · Viết · Cơ bản · A1 · customer
\begin{cauhoi}[ma=jelly-r1-cai-moi-29-8-vn, dulieu=jelly, domain=customer,
  hinhthuc=write, cap=basic, dang=A1, kynang={WHERE, AND, ORDER BY}, bay={quen-dieu-kien-loc, sai-gia-tri-loc},
  tacgia={A}, trangthai=review, bokiemtra={public, hidden-1}, thutu=co, socot=vitri]
\tieude{Ai cài game ở Việt Nam ngày 29/8?}
\boicanh[Chị Mai · Quản lý sản phẩm]{Ngày 29/8 mình chạy thử lời chào mừng mới cho người chơi ở Việt Nam. Em gửi chị danh sách người cài game hôm đó, kèm nền tảng và kênh thu hút.}
\begin{yeucau}
  \item Xét người chơi có ngày cài (\code{install_date}) là 29/8/2026 và thị trường (\code{country}) là Việt Nam.
  \item Trả về \code{user_id}, \code{platform}, \code{channel}.
  \item Sắp \code{user_id} tăng dần.
\end{yeucau}
\begin{erd}
  \erdbang{users}{0,0}{user_id:PK, install_date*=2026-08-14, platform*='Android', country*='VN', channel*='Organic', campaign_id, device_tier}
\end{erd}
\begin{nghiepvu}{Lập danh sách người chơi đã nhận lời chào mừng mới để theo dõi phản hồi.}
  \luuy{Cột \code{country} lưu mã quốc gia hai chữ cái, không lưu tên nước. Mỗi dòng của \code{users} là một người chơi, ghi nhận lúc cài game.}
\end{nghiepvu}
\goiy{Giá trị mẫu của cột \code{country} trong ERD được viết thế nào?}
\goiy{Lọc \code{install_date = '2026-08-29'} và \code{country = 'VN'}.}
\ghichuthietke{Cơ bản: 1 bảng có giá trị mẫu (ngày mẫu tự đặt khác ngày hỏi). Bẫy ở public: quên điều kiện thị trường lấy người cài ở cả 8 thị trường; ghi tên nước \code{'Vietnam'} ra 0 dòng. hidden-1 bỏ người chơi có mã tận cùng bằng 7, nên danh sách khác public, chép kết quả Chạy thử sẽ rớt.}
\end{cauhoi}

% ================================================================ R2 · Viết · Cơ bản · A1 · sales
\begin{cauhoi}[ma=jelly-r2-giao-dich-season-pass, dulieu=jelly, domain=sales,
  hinhthuc=write, cap=basic, dang=A1, kynang={WHERE, ORDER BY}, bay={sai-cot-loc, thieu-can-tren},
  tacgia={A}, trangthai=review, bokiemtra={public, hidden-1}, thutu=co]
\tieude{Ai mua Season Pass trong tháng 7?}
\boicanh[Anh Huy · Quản lý doanh thu]{Đợt Season Pass tháng 7 vừa kết thúc. Anh muốn gửi quà cảm ơn người chơi đã mua gói này, nên cần danh sách giao dịch Season Pass trong tháng.}
\begin{yeucau}
  \item Xét giao dịch gói Season Pass (\code{product_type} là \code{'season_pass'}) trong tháng 7/2026.
  \item Trả về \code{transaction_id}, \code{user_id}, \code{event_time}.
  \item Sắp \code{event_time} tăng dần; cùng thời điểm thì \code{transaction_id} nhỏ hơn đứng trước.
\end{yeucau}
\begin{erd}
  \erdbang{purchases}{0,0}{transaction_id:PK, user_id:FK, event_date*=2026-07-20, event_time*=2026-07-20 19:42:05, product_type*='season_pass', product_name='Season Pass', price_usd=9.99}
\end{erd}
\begin{nghiepvu}{Lập danh sách người chơi nhận quà cảm ơn khi đợt Season Pass tháng 7 kết thúc.}
  \khainiem{doanh-thu-iap}
  \luuy{\code{product_type} là mã gói, viết thường và nối bằng gạch dưới; \code{product_name} là tên gói hiện trong cửa hàng của game. Hai cột ghi cùng một gói theo hai cách khác nhau.}
\end{nghiepvu}
\goiy{So giá trị mẫu của \code{product_type} và \code{product_name} trong ERD. Giá trị \code{'season_pass'} thuộc cột nào?}
\goiy{Lọc \code{product_type = 'season_pass'} và \code{event_date} từ \code{'2026-07-01'} tới trước \code{'2026-08-01'}.}
\ghichuthietke{Cơ bản: 1 bảng có giá trị mẫu. Bẫy ở public: lọc \code{product_name = 'season_pass'} ra 0 dòng (collation không phân biệt hoa thường nhưng dấu gạch dưới khác dấu cách); thiếu cận trên lấy cả giao dịch Season Pass tháng 8. hidden-1 bỏ khoảng 10\% người chơi nên danh sách khác public.}
\end{cauhoi}

% ================================================================ R3 · Viết · Cơ bản · A1 · operations
\begin{cauhoi}[ma=jelly-r3-man-sieu-kho, dulieu=jelly, domain=operations,
  hinhthuc=write, cap=basic, dang=A1, kynang={WHERE, ORDER BY}, bay={like-rong, sai-gia-tri-loc},
  tacgia={A}, trangthai=review, bokiemtra={public, hidden-1}, thutu=co]
\tieude{Màn nào được gắn nhãn siêu khó?}
\boicanh[Anh Khoa · Thiết kế màn chơi]{Trước khi thiết kế thế giới mới, anh muốn xem lại các màn khó nhất của game. Em liệt kê giúp anh những màn có nhãn siêu khó, kèm thế giới và số lượt đi tối đa.}
\begin{yeucau}
  \item Xét các màn có nhãn độ khó (\code{difficulty_tag}) là siêu khó.
  \item Trả về \code{level_id}, \code{world_name}, \code{moves_limit}.
  \item Sắp \code{level_id} tăng dần.
\end{yeucau}
\begin{erd}
  \erdbang{level_metadata}{0,0}{level_id:PK, episode, world_name*='Jelly Jungle', difficulty_tag*='hard', moves_limit*=29, designed_fail_probability}
\end{erd}
\begin{nghiepvu}{Xác định màn cần xem lại số lượt đi và mục tiêu trước khi thiết kế thế giới mới.}
  \luuy{Nhãn độ khó có ba mức: \code{'normal'} (thường), \code{'hard'} (khó) và \code{'super_hard'} (siêu khó). Mỗi màn chỉ mang một nhãn.}
\end{nghiepvu}
\goiy{So khớp một phần chuỗi theo chữ \code{hard} có lấy cả màn mang nhãn \code{'hard'} không?}
\goiy{So bằng đúng giá trị: \code{difficulty_tag = 'super_hard'}.}
\ghichuthietke{Cơ bản: 1 bảng có giá trị mẫu; giá trị mẫu \code{'hard'} là manh mối của bẫy. Đề nói ``siêu khó'' bằng lời, học viên tự tra giá trị trong tab Nghiệp vụ. Bẫy ở public: \code{LIKE} có ký tự đại diện hai đầu lấy cả 4 màn \code{'hard'}; lọc khác \code{'normal'} cũng lấy cả hai nhãn. hidden-1 đổi nhãn màn 45 từ \code{'hard'} thành \code{'super_hard'} (5 dòng thay vì 4), nên bài chép 4 mã màn từ lần Chạy thử rớt.}
\end{cauhoi}

% ================================================================ R4 · Viết · Dễ · A1 · marketing
\begin{cauhoi}[ma=jelly-r4-ngay-chi-cao-social, dulieu=jelly, domain=marketing,
  hinhthuc=write, cap=easy, dang=A1, kynang={WHERE, ORDER BY, LIMIT}, bay={quen-dieu-kien-loc, thieu-can-tren},
  tacgia={A}, trangthai=review, bokiemtra={public, hidden-1}, thutu=co, saiso=0.01]
\tieude{Mười ngày chi nhiều nhất cho Social Ads ở Việt Nam}
\boicanh[Anh Tuấn · Quản lý thu hút người chơi]{Từ ngày 15/8 anh tăng giá thầu cho chiến dịch quảng cáo trên mạng xã hội nhắm tới người chơi Android ở Việt Nam. Em lấy giúp anh mười ngày chiến dịch này tốn nhiều tiền nhất, tính từ hôm đó tới nay.}
\begin{yeucau}
  \item Xét chi phí của chiến dịch \code{'JTA-SOCIAL-204'} (kênh Social Ads), thị trường \code{'VN'}, nền tảng nhắm tới (\code{target_platform}) \code{'Android'}, từ ngày 15/8/2026 tới ngày mới nhất có trong dữ liệu.
  \item Trả về \code{date}, \code{spend_usd} (USD).
  \item Sắp \code{spend_usd} giảm dần; bằng nhau thì ngày sớm hơn đứng trước.
  \item Lấy 10 dòng đầu.
\end{yeucau}
\begin{erd}
  \erdbang{campaign_spend}{0,0}{date*, campaign_id*, channel, country*, target_platform*, spend_usd*}
  \erdbang{users}{1,0}{user_id:PK, install_date, channel, campaign_id}
  \erdnoi[tren={nối logic theo chiến dịch, nhiều -- nhiều}]{campaign_spend.campaign_id}{users.campaign_id}
\end{erd}
\begin{nghiepvu}{Xác định ngày cần xem lại giá thầu sau khi tăng giá thầu chiến dịch.}
  \luuy{Mỗi dòng của \code{campaign_spend} là chi phí của một chiến dịch trong một ngày, ở một thị trường, trên một nền tảng nhắm tới. Một ngày của cùng chiến dịch ở Việt Nam có thể có hai dòng, một cho iOS và một cho Android.}
\end{nghiepvu}
\goiy{Chiến dịch này nhắm tới cả iOS lẫn Android. Câu hỏi chỉ cần nền tảng nào?}
\goiy{Thêm \code{target_platform = 'Android'} và chỉ đặt cận dưới \code{date >= '2026-08-15'}; ngày cuối của lần Chạy thử chưa chắc là ngày cuối của dữ liệu chấm.}
\ghichuthietke{Dễ: 2 bảng, chữ trên đường nối (nối logic, không có khoá ngoại); câu chỉ cần \code{campaign_spend}. v2: đổi phạm vi thành ``từ 15/8 tới ngày mới nhất'' để kết quả đúng của hidden-1 khác public (biến thể có thêm dòng chi ngày 2026-09-01 của đúng chiến dịch, thị trường, nền tảng này, đứng hạng 3), nên bài chép kết quả Chạy thử rớt. Bẫy ở public: thiếu điều kiện nền tảng thì dòng iOS ngày 23/8 lọt vào 10 dòng đầu. Bẫy ở hidden-1: gắn cứng ngày cuối \code{'2026-08-29'} của dữ liệu Chạy thử làm mất ngày chi mới.}
\end{cauhoi}

% ================================================================ R5 · Viết · Dễ · A1 · operations
\begin{cauhoi}[ma=jelly-r5-luot-dai-man-30, dulieu=jelly, domain=operations,
  hinhthuc=write, cap=easy, dang=A1, kynang={WHERE, ORDER BY, LIMIT}, bay={phien-ban-luc-cai, quen-dieu-kien-loc},
  tacgia={A}, trangthai=review, bokiemtra={public, hidden-1}, thutu=co]
\tieude{Lượt chơi màn 30 lâu nhất trên Android 2.1.0}
\boicanh[Chị Linh · Kỹ sư QA]{Sau bản 2.1.0, một số người chơi Android báo màn 30 chạy giật và chơi rất lâu mới xong. Em lấy giúp chị hai mươi lượt chơi lâu nhất trên bản này để nhóm tái hiện lỗi.}
\begin{yeucau}
  \item Xét lượt chơi màn 30 (\code{level_id}) trên nền tảng \code{'Android'}, với phiên bản ứng dụng lúc chơi (\code{app_version}) là \code{'2.1.0'}.
  \item Trả về \code{attempt_id}, \code{user_id}, \code{event_date}, \code{duration_sec} (giây), \code{result}.
  \item Sắp \code{duration_sec} giảm dần; bằng nhau thì \code{attempt_id} nhỏ hơn đứng trước.
  \item Lấy 20 dòng đầu.
\end{yeucau}
\begin{erd}
  \erdbang{level_attempts}{0,0}{attempt_id:PK, user_id:FK, event_date*, level_id*=30, result*, duration_sec*, app_version*='2.1.1', platform*='Android'}
  \erdbang{users}{1,0}{user_id:PK, install_date, platform, app_version_install='2.0.2'}
  \erdnoi[kieu=n-1, tren={\code{level_attempts.user_id = users.user_id}}]{level_attempts.user_id}{users.user_id}
\end{erd}
\begin{nghiepvu}{Lấy lượt chơi để nhóm QA tái hiện lỗi trên đúng phiên bản.}
  \khainiem{phien-ban-luc-choi}
  \luuy{Người chơi có thể cập nhật ứng dụng sau khi cài, nên phiên bản lúc cài và phiên bản của một lượt chơi có thể khác nhau.}
\end{nghiepvu}
\goiy{Phiên bản người chơi dùng khi chơi lượt đó được ghi ở bảng nào?}
\goiy{Lọc \code{app_version} của \code{level_attempts}; câu này không cần nối với \code{users}.}
\ghichuthietke{Dễ: 2 bảng, chữ trên đường nối; \code{users} có mặt để làm rõ hai cột phiên bản. Có nhiều lượt bằng thời lượng quanh hạng 20 nên đề ghi tiêu chí phụ. Bẫy ở public: lọc theo \code{users.app_version_install} (phiên bản lúc cài) ra danh sách khác; quên điều kiện nền tảng lẫn lượt iOS. hidden-1 bỏ khoảng 10\% lượt chơi nên danh sách khác public.}
\end{cauhoi}

% ================================================================ R6 · Viết · Cơ bản · B4 · marketing
\begin{cauhoi}[ma=jelly-r6-su-kien-thang-8, dulieu=jelly, domain=marketing,
  hinhthuc=write, cap=basic, dang=B4, kynang={WHERE, AND/OR, so sánh ngày}, bay={hoac-thay-va, giao-khoang-ngay},
  tacgia={A}, trangthai=review, bokiemtra={public, hidden-1}, thutu=co]
\tieude{Sự kiện LiveOps nào chạy từ 26/7 tới 9/8?}
\boicanh[Chị Thảo · Trưởng nhóm LiveOps]{Chị đang viết báo cáo vận hành hai tuần từ 26/7 tới 9/8. Em liệt kê giúp chị các sự kiện trong game có ít nhất một ngày chạy trong hai tuần đó.}
\begin{yeucau}
  \item Xét các sự kiện LiveOps (sự kiện vận hành trực tiếp trong game) có ít nhất một ngày chạy trong khoảng từ 26/7/2026 tới hết 9/8/2026, dựa vào ngày bắt đầu (\code{start_date}) và ngày kết thúc (\code{end_date}).
  \item Trả về \code{event_name}, \code{start_date}, \code{end_date}.
  \item Sắp \code{start_date} tăng dần.
\end{yeucau}
\begin{erd}
  \erdbang{liveops_events}{0,0}{event_id:PK, event_name*='Weekend Sugar Rush', start_date*=2026-07-10, end_date*=2026-07-12, mechanic}
\end{erd}
\begin{nghiepvu}{Lập danh sách sự kiện đưa vào báo cáo vận hành hai tuần.}
  \luuy{Sự kiện LiveOps chạy theo một khoảng ngày. Sự kiện bắt đầu trước kỳ báo cáo, hoặc kết thúc sau kỳ báo cáo, vẫn có ngày chạy trong kỳ.}
\end{nghiepvu}
\goiy{Sự kiện chạy từ 24/7 tới 28/7 có ngày nào nằm trong kỳ không? Ngày bắt đầu của sự kiện đó có nằm trong kỳ không?}
\goiy{Hai khoảng ngày giao nhau khi sự kiện bắt đầu không muộn hơn ngày cuối kỳ \code{AND} kết thúc không sớm hơn ngày đầu kỳ.}
\ghichuthietke{Cơ bản: 1 bảng có giá trị mẫu (sự kiện có thật ngoài kỳ). Dạng B4 vì hỏi sự kiện đang chạy trong một khoảng thời gian. v2: đổi kỳ từ tháng 8 sang 26/7 -- 9/8 để có hai sự kiện vắt qua mốc đầu (24/7 -- 28/7) và mốc cuối (8/8 -- 12/8); mã câu giữ nguyên. Bẫy ở public: chỉ xét ngày bắt đầu trong kỳ ra 1 dòng; đòi nằm trọn trong kỳ ra 0 dòng; nối bằng \code{OR} ra 3 dòng. hidden-1 thêm sự kiện Jungle Jam Weekend (31/7 -- 2/8) nằm trọn trong kỳ (3 dòng thay vì 2), nên bài chép kết quả Chạy thử rớt.}
\end{cauhoi}

% ================================================================ R7 · Sửa lỗi · Cơ bản · A1 · customer
\begin{cauhoi}[ma=jelly-r7-organic-ios-tuan-dau, dulieu=jelly, domain=customer,
  hinhthuc=debug, cap=basic, dang=A1, kynang={AND/OR, ngoặc, IN}, bay={thu-tu-and-or},
  tacgia={A}, trangthai=review, bokiemtra={public, hidden-1}, thutu=co]
\tieude{Danh sách người chơi iOS bị lẫn người chơi Android}
\boicanh[Bạn Nam · Thực tập sinh phân tích]{Em lấy danh sách người chơi iOS ở Việt Nam cài tuần đầu tháng 7, tự tìm thấy game hoặc đến qua quảng cáo tìm kiếm. Danh sách lại có cả người chơi Android và người cài tháng 8, em chưa nhìn ra mình lọc nhầm chỗ nào.}
\begin{yeucau}
  \item Sửa truy vấn dưới đây để chỉ lấy người chơi ở thị trường \code{'VN'}, nền tảng \code{'iOS'}, cài từ ngày 1/7/2026 tới hết ngày 7/7/2026, đến từ kênh \code{'Organic'} hoặc \code{'Search Ads'}.
  \item Giữ nguyên các cột \code{user_id}, \code{install_date}, \code{channel} và thứ tự sắp.
\end{yeucau}
\begin{erd}
  \erdbang{users}{0,0}{user_id:PK, install_date*=2026-07-03, platform*='iOS', country*='VN', channel*='Search Ads', campaign_id}
\end{erd}
\begin{nghiepvu}{Lập nhóm người chơi iOS để so cách người chơi tìm thấy game.}
  \luuy{Kênh \code{'Organic'} là người chơi tự tìm thấy game, không qua quảng cáo trả phí; \code{'Search Ads'} là quảng cáo trên kết quả tìm kiếm của cửa hàng ứng dụng.}
\end{nghiepvu}
\goiy{Trong điều kiện lọc, \code{AND} được tính trước \code{OR}. Điều kiện \code{channel = 'Search Ads'} đang được ghép với các điều kiện còn lại như thế nào?}
\ghichuthietke{Cơ bản: lỗi một chỗ (thiếu ngoặc quanh hai điều kiện kênh). Starter lấy thêm mọi người chơi Search Ads ở mọi thị trường, nền tảng, ngày cài. Bài sai ở public: bỏ hẳn điều kiện Search Ads (mất một phần danh sách); đổi \code{OR} thành \code{AND} ra 0 dòng. Lời giải dùng ngoặc; \code{channel IN ('Organic', 'Search Ads')} cũng đúng.}
\end{cauhoi}

% ================================================================ R8 · Sửa lỗi · Dễ · A1 · sales
\begin{cauhoi}[ma=jelly-r8-giao-dich-thang-8-thua, dulieu=jelly, domain=sales,
  hinhthuc=debug, cap=easy, dang=A1, kynang={WHERE, khoảng ngày nửa mở}, bay={thieu-can-tren, between-gom-dau-mut, lon-hon-thay-lon-hon-bang},
  tacgia={A}, trangthai=review, bokiemtra={public, hidden-1}, thutu=co]
\tieude{Danh sách giao dịch tháng 8 có dòng tháng 9}
\boicanh[Chị Hà · Kế toán quản trị]{Mỗi sáng chị chạy lại báo cáo giao dịch tháng 8 của thị trường Anh. Từ khi có thêm dữ liệu ngày mới, danh sách lẫn cả giao dịch tháng 9. Em sửa giúp chị.}
\begin{yeucau}
  \item Danh sách chỉ được gồm giao dịch của người chơi ở thị trường \code{'GB'} trong tháng 8/2026; sửa truy vấn dưới đây cho đúng.
  \item Truy vấn vẫn phải đúng khi dữ liệu có thêm ngày sau tháng 8.
  \item Giữ nguyên các cột \code{transaction_id}, \code{user_id}, \code{event_date}, \code{product_type}, \code{price_usd} và thứ tự sắp.
\end{yeucau}
\begin{erd}
  \erdbang{purchases}{0,0}{transaction_id:PK, user_id:FK, event_date*=2026-09-01, product_type*='booster_bundle', price_usd*=4.99, country*='GB'}
  \erdbang{users}{1,0}{user_id:PK, install_date, country}
  \erdnoi[kieu=n-1, tren={\code{purchases.user_id = users.user_id}}]{purchases.user_id}{users.user_id}
\end{erd}
\begin{nghiepvu}{Đối chiếu doanh thu tháng 8 của thị trường Anh với số liệu cửa hàng ứng dụng gửi về.}
  \khainiem{doanh-thu-iap}
  \luuy{Dữ liệu dùng khi Chạy thử chỉ tới ngày 29/8/2026, còn báo cáo thật chạy trên dữ liệu có thêm ngày mới. Điều kiện lọc theo tháng cần cả mốc đầu và mốc cuối.}
\end{nghiepvu}
\goiy{Một giao dịch ngày 1/9/2026 (như giá trị mẫu trong ERD) có thoả điều kiện ngày hiện tại không?}
\goiy{Thêm cận trên \code{event_date < '2026-09-01'}, giữ cận dưới \code{>= '2026-08-01'}.}
\ghichuthietke{Dễ: 2 bảng, chữ trên đường nối; giá trị mẫu 2026-09-01 là manh mối. Starter thiếu cận trên: khớp lời giải ở public (dữ liệu dừng 29/8) nhưng rớt ở hidden-1, vì biến thể có thêm một giao dịch Booster Bundle ngày 2026-09-01 của người chơi ở GB. v2: bỏ điều kiện gói để có giao dịch ngày 1/8 ở public. Bài sai: sửa thành \code{BETWEEN '2026-08-01' AND '2026-09-01'} vẫn gồm ngày 1/9 (rớt ở hidden-1); thêm cận trên nhưng đổi cận dưới thành \code{>} mất các giao dịch ngày 1/8 (rớt ở public). Cách gắn cứng \code{<= '2026-08-29'} chưa bắt được vì không biến thể nào có giao dịch 30/8 hoặc 31/8; chờ giảng viên.}
\end{cauhoi}

% ================================================================ R9 · Sửa lỗi · Dễ · A1 · operations
\begin{cauhoi}[ma=jelly-r9-phien-dai-15-8, dulieu=jelly, domain=operations,
  hinhthuc=debug, cap=easy, dang=A1, kynang={ORDER BY DESC, tiêu chí phụ, LIMIT}, bay={chieu-sap-xep, thieu-tieu-chi-phu},
  tacgia={A}, trangthai=review, bokiemtra={public, hidden-1}, thutu=co]
\tieude{Năm phiên dài nhất ngày 15/8 chỉ dài vài phút}
\boicanh[Bạn Vy · Thực tập sinh phân tích]{Em lấy năm phiên chơi dài nhất ngày 15/8 để gửi nhóm thiết kế, nhưng các phiên trong danh sách chỉ dài vài phút. Chiều nay em phải gửi rồi, nhờ anh chị sửa giúp.}
\begin{yeucau}
  \item Truy vấn dưới đây cần lấy 5 phiên chơi có thời lượng (\code{duration_minutes}) dài nhất trong ngày 15/8/2026 nhưng đang lấy sai; sửa lại.
  \item Hai phiên bằng thời lượng thì \code{session_id} nhỏ hơn đứng trước.
  \item Giữ nguyên các cột \code{session_id}, \code{user_id}, \code{session_start}, \code{duration_minutes}.
\end{yeucau}
\begin{erd}
  \erdbang{sessions}{0,0}{session_id:PK, user_id:FK, session_start*=2026-08-15 20:31:00, event_date*=2026-08-15, duration_minutes*=12, entry_screen}
  \erdbang{users}{1,0}{user_id:PK, platform, country}
  \erdnoi[kieu=n-1, tren={\code{sessions.user_id = users.user_id}}]{sessions.user_id}{users.user_id}
\end{erd}
\begin{nghiepvu}{Gửi nhóm thiết kế các phiên chơi dài để xem người chơi làm gì trong phiên.}
  \khainiem{phien-choi}
  \luuy{Một người chơi có thể có nhiều phiên trong một ngày. \code{session_start} ghi cả ngày và giờ mở phiên, \code{event_date} chỉ ghi ngày.}
\end{nghiepvu}
\goiy{Truy vấn đang sắp thời lượng theo chiều nào? Năm dòng đầu của chiều đó là phiên dài hay phiên ngắn?}
\goiy{Sắp \code{duration_minutes} giảm dần rồi \code{session_id} tăng dần trước khi lấy 5 dòng.}
\ghichuthietke{Dễ: 2 bảng, chữ trên đường nối. Starter có hai lỗi: sắp tăng dần và thiếu tiêu chí phụ (hạng 5 có hai phiên bằng thời lượng ở public). Bài sai: sắp tiêu chí phụ ngược chiều (\code{session_id DESC}) làm hai phiên 41 phút đổi chỗ (rớt ở public); so \code{session_start = '2026-08-15'} chỉ khớp phiên mở đúng 0 giờ (rớt ở public). hidden-1 bỏ khoảng 10\% phiên nên danh sách khác public.}
\end{cauhoi}


% ################################################################ MODULE 2 · NULL & EXPRESSIONS

% ================================================================ N1 · Viết · Cơ bản · A2 · sales
\begin{cauhoi}[ma=jelly-n1-iap-khong-gan-man, dulieu=jelly, domain=sales,
  hinhthuc=write, cap=basic, dang=A2, kynang={IS NULL, COUNT}, bay={bang-null, count-cot-null},
  tacgia={A}, trangthai=review, bokiemtra={public, hidden-1}, thutu=co]
\tieude{Bao nhiêu giao dịch không gắn với màn chơi?}
\boicanh[Anh Huy · Quản lý doanh thu]{Anh muốn biết người chơi thường mua ngoài màn chơi hay mua ngay lúc đang chơi một màn. Trước hết em đếm giúp anh số giao dịch không gắn với màn nào.}
\begin{yeucau}
  \item Xét mọi giao dịch trong bảng \code{purchases}.
  \item Đếm các giao dịch có \code{level_id_if_relevant} trống (NULL).
  \item Trả về một cột \code{so_giao_dich}.
\end{yeucau}
\begin{erd}
  \erdbang{purchases}{0,0}{transaction_id:PK, product_type, level_id_if_relevant*=NULL, screen_context='Shop screen', price_usd}
\end{erd}
\begin{nghiepvu}{Quyết định đặt ưu đãi mua hàng trong cửa hàng hay ngay trong màn chơi.}
  \khainiem{doanh-thu-iap}
  \luuy{Giao dịch không gắn với màn chơi nào có \code{level_id_if_relevant} để trống. Một số gói mua ở cửa hàng vẫn ghi màn chơi đi kèm, nên ô trống không trùng với giao dịch ở cửa hàng. Ô trống (NULL) nghĩa là không có giá trị, khác với số 0.}
\end{nghiepvu}
\goiy{Phép so bằng (\code{=}) giữa một ô trống và \code{NULL} cho ra kết quả gì?}
\goiy{Lọc bằng \code{IS NULL} rồi đếm số dòng bằng \code{COUNT(*)}.}
\ghichuthietke{Cơ bản: 1 bảng, giá trị mẫu NULL là manh mối. Bẫy ở public: \code{= NULL} ra 0; \code{COUNT(level_id_if_relevant)} bỏ qua ô trống nên cũng ra 0. hidden-1 bỏ khoảng 10\% người chơi nên số khác public.}
\end{cauhoi}

% ================================================================ N2 · Viết · Cơ bản · A1 · marketing
\begin{cauhoi}[ma=jelly-n2-chua-tiep-xuc-ab, dulieu=jelly, domain=marketing,
  hinhthuc=write, cap=basic, dang=A1, kynang={IS NULL, ORDER BY, LIMIT}, bay={bang-null, dao-dieu-kien},
  tacgia={A}, trangthai=review, bokiemtra={public, hidden-1}, thutu=co]
\tieude{Ai được chia nhánh nhưng chưa tiếp xúc ưu đãi?}
\boicanh[Chị Mai · Quản lý sản phẩm]{Thử nghiệm ưu đãi thêm lượt đi đã chia nhánh xong, nhưng có người chơi chưa thấy ưu đãi. Em lấy giúp chị hai mươi người đầu tiên theo mã người chơi để nhóm QA kiểm tra.}
\begin{yeucau}
  \item Xét thử nghiệm A/B \code{'extra_moves_offer_timing'}.
  \item Lọc người chơi chưa tiếp xúc ưu đãi (\code{exposed_at} trống).
  \item Trả về \code{user_id}, \code{variant}.
  \item Sắp \code{user_id} tăng dần; lấy 20 dòng đầu.
\end{yeucau}
\begin{erd}
  \erdbang{ab_assignments}{0,0}{user_id:PK, experiment_id:PK*, variant*='control_after_3_fails', assigned_at, exposed_at*=NULL, exposure_status='not_exposed'}
\end{erd}
\begin{nghiepvu}{Kiểm tra vì sao một phần người chơi được chia nhánh chưa tiếp xúc ưu đãi.}
  \khainiem{thu-nghiem-ab}
  \luuy{Người chơi chỉ tiếp xúc ưu đãi khi thua màn theo đúng quy tắc của nhánh mình. Người chưa tiếp xúc ưu đãi vẫn thuộc nhánh được chia.}
\end{nghiepvu}
\goiy{Điều kiện \code{exposed_at = NULL} đúng với dòng nào?}
\goiy{Dùng \code{exposed_at IS NULL}.}
\ghichuthietke{Cơ bản: 1 bảng, giá trị mẫu NULL là manh mối. Bẫy ở public: \code{= NULL} ra 0 dòng; \code{IS NOT NULL} lấy người đã tiếp xúc. Lọc \code{exposure_status = 'not_exposed'} cho cùng kết quả (hai cột khớp nhau), chấp nhận. hidden-1 bỏ người chơi có mã tận cùng bằng 7 nên 20 dòng đầu khác public.}
\end{cauhoi}

% ================================================================ N3 · Viết · Dễ · A3 · operations
\begin{cauhoi}[ma=jelly-n3-nhom-buoc-du, dulieu=jelly, domain=operations,
  hinhthuc=write, cap=easy, dang=A3, kynang={CASE WHEN}, bay={thu-tu-case, nham-ket-qua-voi-luot-du},
  tacgia={A}, trangthai=review, bokiemtra={public, hidden-1}, thutu=co]
\tieude{Lượt chơi màn 28 còn dư bao nhiêu lượt đi?}
\boicanh[Anh Khoa · Thiết kế màn chơi]{Anh nghi màn 28 cho quá ít lượt đi. Em chia giúp anh các lượt chơi màn này ngày 1/8 thành nhóm theo số lượt đi còn dư lúc kết thúc.}
\begin{yeucau}
  \item Xét các lượt chơi màn 28 (\code{level_id}) ngày 1/8/2026.
  \item Gán nhóm theo số lượt đi còn dư (\code{moves_left}): bằng 0 là \code{'het_luot_di'}; từ 1 tới 2 là \code{'du_1_2'}; từ 3 trở lên là \code{'du_tu_3'}.
  \item Trả về \code{attempt_id}, \code{result}, \code{moves_left}, \code{nhom_luot_du}.
  \item Sắp \code{attempt_id} tăng dần.
\end{yeucau}
\begin{erd}
  \erdbang{level_attempts}{0,0}{attempt_id:PK, event_date*, level_id:FK*, result*, moves_limit, moves_used, moves_left*}
  \erdbang{level_metadata}{1,0}{level_id:PK, world_name, difficulty_tag, moves_limit}
  \erdnoi[kieu=n-1, tren={\code{level_attempts.level_id = level_metadata.level_id}}]{level_attempts.level_id}{level_metadata.level_id}
\end{erd}
\begin{nghiepvu}{Quyết định có tăng số lượt đi tối đa của màn hay không.}
  \luuy{Số lượt đi đã dùng cộng số lượt đi còn dư bằng số lượt đi tối đa của màn. Lượt thắng vẫn có thể còn 0 lượt đi khi người chơi đạt mục tiêu đúng ở lượt đi cuối. Lượt bỏ cuộc và lượt gặp sự cố được ghi còn 0 lượt đi.}
\end{nghiepvu}
\goiy{Một lượt có \code{moves_left} bằng 0 thoả những nhánh điều kiện nào trong ba nhóm?}
\goiy{\code{CASE WHEN} dừng ở nhánh đúng đầu tiên; đặt nhánh \code{moves_left = 0} lên trước nhánh \code{moves_left <= 2}.}
\ghichuthietke{Dễ: 2 bảng, chữ trên đường nối; câu chỉ cần \code{level_attempts}. Nhãn viết không dấu, có gạch dưới để tránh lệch chuẩn Unicode khi so chuỗi. Bẫy ở public: đặt nhánh \code{<= 2} trước nên lượt còn 0 lượt đi bị gán \code{'du_1_2'}; gán theo kết quả (\code{result}) thay vì số lượt dư làm sai các lượt thắng còn 0 lượt đi. Màn 28 thuộc vùng khó đột ngột 28--30. hidden-1 bỏ khoảng 10\% lượt chơi nên danh sách khác public.}
\end{cauhoi}

% ================================================================ N4 · Viết · Dễ · B1 · finance
\begin{cauhoi}[ma=jelly-n4-tong-doanh-thu-ngay, dulieu=jelly, domain=finance,
  hinhthuc=write, cap=easy, dang=B1, kynang={phép tính, ROUND, thứ tự phép tính}, bay={thu-tu-phep-tinh, sai-mau-so},
  tacgia={A}, trangthai=review, bokiemtra={public, hidden-1}, thutu=co, saiso=0.01]
\tieude{Tổng doanh thu từng ngày tuần đầu tháng 7}
\boicanh[Chị Ngọc · Quản lý doanh thu]{Chị đang xem lại tuần đầu tiên sau khi game ra mắt. Em cho chị biết mỗi ngày game thu được bao nhiêu và quảng cáo chiếm bao nhiêu phần trong đó.}
\begin{yeucau}
  \item Xét 7 ngày từ 1/7/2026 tới hết 7/7/2026 trong bảng chỉ số ngày \code{daily_kpi}.
  \item Tính tổng doanh thu bằng doanh thu IAP (\code{iap_revenue}) cộng doanh thu quảng cáo (\code{ad_revenue}); tính tỷ trọng quảng cáo bằng doanh thu quảng cáo chia tổng doanh thu, đổi ra phần trăm.
  \item Trả về \code{date}, \code{tong_doanh_thu} (USD), \code{ty_trong_quang_cao_pct} (phần trăm, làm tròn 2 chữ số).
  \item Sắp \code{date} tăng dần.
\end{yeucau}
\begin{erd}
  \erdbang{daily_kpi}{0,0}{date:PK*, dau, iap_revenue*=40.15, ad_revenue*=5.12, arpdau}
  \erdbang{player_daily}{1,0}{user_id:PK, activity_date:PK, iap_revenue, ad_revenue}
  \erdnoi[kieu=1-n, tren={tổng hợp theo ngày, không có khoá}]{daily_kpi.date}{player_daily.activity_date}
\end{erd}
\begin{nghiepvu}{Theo dõi phần doanh thu đến từ quảng cáo trong những ngày đầu ra mắt.}
  \khainiem{doanh-thu-iap}
  \khainiem{doanh-thu-quang-cao}
  \luuy{Tỷ trọng quảng cáo là phần doanh thu quảng cáo trong tổng doanh thu, nên mẫu số là tổng doanh thu IAP và quảng cáo. Bảng chỉ số ngày (\code{daily_kpi}) được tổng hợp sẵn từ \code{player_daily}, mỗi dòng là một ngày.}
\end{nghiepvu}
\goiy{Trong biểu thức có cả phép cộng và phép chia, phép nào được làm trước?}
\goiy{Đặt \code{iap_revenue + ad_revenue} trong ngoặc khi dùng làm mẫu số, rồi nhân 100 và làm tròn 2 chữ số.}
\ghichuthietke{Dễ: 2 bảng, đường nối ghi rõ là tổng hợp theo ngày; câu chỉ cần \code{daily_kpi}. Bản đồ ghi ``ARPDAU tự tính''; đổi sang tỷ trọng quảng cáo để không trùng câu n8 (ARPDAU), và gắn dạng B1 vì câu theo dõi một chỉ số qua chuỗi 7 ngày. Bẫy ở public: thiếu ngoặc mẫu số; chia cho doanh thu IAP. hidden-1 tính lại \code{daily_kpi} sau khi bỏ khoảng 10\% người chơi nên số khác public.}
\end{cauhoi}

% ================================================================ N5 · Viết · Dễ · A1 · sales
\begin{cauhoi}[ma=jelly-n5-gia-quy-doi-vnd, dulieu=jelly, domain=sales,
  hinhthuc=write, cap=easy, dang=A1, kynang={DISTINCT, phép tính, ROUND âm}, bay={thieu-distinct, lam-tron-sai},
  tacgia={A}, trangthai=review, bokiemtra={public, hidden-1}, thutu=co, saiso=0.01]
\tieude{Giá các gói vật phẩm quy ra đồng}
\boicanh[Chị Ngọc · Quản lý doanh thu]{Chị chuẩn bị trình bảng giá bán ở thị trường Việt Nam. Em lập giúp chị giá bán từng gói vật phẩm quy ra đồng, làm tròn tới nghìn đồng.}
\begin{yeucau}
  \item Lấy mỗi cặp gói vật phẩm – mức giá đã bán một dòng từ bảng giao dịch \code{purchases}.
  \item Quy đổi giá theo tỷ giá 25.400 đồng một USD, làm tròn tới hàng nghìn.
  \item Trả về \code{product_name}, \code{price_usd} (USD), \code{gia_vnd} (đồng).
  \item Sắp \code{price_usd} tăng dần.
\end{yeucau}
\begin{erd}
  \erdbang{purchases}{0,0}{transaction_id:PK, product_type, product_name*='Piggy Bank', price_usd*=2.99, level_id_if_relevant:FK}
  \erdbang{level_metadata}{1,0}{level_id:PK, world_name}
  \erdnoi[kieu=n-1, tren={\code{purchases.level_id_if_relevant = level_metadata.level_id}}]{purchases.level_id_if_relevant}{level_metadata.level_id}
\end{erd}
\begin{nghiepvu}{Chuẩn bị bảng giá bán quy ra đồng để trình lãnh đạo duyệt.}
  \luuy{Mỗi dòng của \code{purchases} là một giao dịch, nên cùng một gói xuất hiện ở nhiều dòng. Gói nào từng bán ở hơn một mức giá thì mỗi mức giá là một dòng.}
\end{nghiepvu}
\goiy{Bảng giao dịch có bao nhiêu dòng cho cùng một gói? Làm sao chỉ giữ mỗi gói một dòng?}
\goiy{Dùng \code{SELECT DISTINCT}; \code{ROUND} với số chữ số âm làm tròn phần nguyên, \code{-3} là làm tròn tới hàng nghìn.}
\ghichuthietke{Dễ: 2 bảng, chữ trên đường nối; câu chỉ cần \code{purchases}. Tỷ giá đặt trong đề, không lấy từ dữ liệu. Bẫy ở public: thiếu \code{DISTINCT} ra hàng nghìn dòng; \code{ROUND(x, 3)} giữ phần lẻ thay vì làm tròn tới nghìn. Hai bài sai rớt cả hai biến thể. hidden-1 có thêm giao dịch Starter Pack giá khuyến mãi 0,49 USD, nên solution ra 6 dòng (Starter Pack hai mức giá) và bài chép bảng giá từ lần Chạy thử rớt; yêu cầu ghi ``mỗi cặp gói – mức giá một dòng'' để khớp cả hai biến thể.}
\end{cauhoi}

% ================================================================ N6 · Viết · Dễ · A2 · finance
\begin{cauhoi}[ma=jelly-n6-ngay-choi-khong-doanh-thu, dulieu=jelly, domain=finance,
  hinhthuc=write, cap=easy, dang=A2, kynang={COUNT(CASE WHEN), phép tính}, bay={lam-tron-truoc-khi-so, thieu-can-tren},
  tacgia={A}, trangthai=review, bokiemtra={public, hidden-1}, thutu=co]
\tieude{Bao nhiêu ngày chơi không phát sinh doanh thu?}
\boicanh[Anh Phong · Quản lý doanh thu]{Anh đang cân nhắc đặt thêm quảng cáo trong game. Em cho anh biết tháng 7 có bao nhiêu ngày chơi không mang lại đồng doanh thu nào, kể cả doanh thu quảng cáo.}
\begin{yeucau}
  \item Xét các dòng \code{player_daily} có ngày hoạt động (\code{activity_date}) trong tháng 7/2026; mỗi dòng là một ngày chơi của một người chơi.
  \item Đếm tổng số ngày chơi và số ngày chơi có tổng doanh thu IAP (\code{iap_revenue}) và doanh thu quảng cáo (\code{ad_revenue}) bằng 0.
  \item Trả về \code{so_ngay_choi}, \code{so_ngay_khong_doanh_thu}.
\end{yeucau}
\begin{erd}
  \erdbang{player_daily}{0,0}{user_id:PK, activity_date:PK*, sessions, iap_revenue*=0.00, ad_revenue*=0.0031}
  \erdbang{users}{1,0}{user_id:PK, install_date, country}
  \erdnoi[kieu=n-1, tren={\code{player_daily.user_id = users.user_id}}]{player_daily.user_id}{users.user_id}
\end{erd}
\begin{nghiepvu}{Ước lượng phần ngày chơi chưa tạo ra doanh thu trước khi cân nhắc thêm vị trí đặt quảng cáo.}
  \khainiem{doanh-thu-quang-cao}
  \luuy{Bảng \code{player_daily} chỉ có ngày người chơi có mở game. Doanh thu quảng cáo của một ngày có thể nhỏ hơn 0,01 USD nhưng vẫn khác 0.}
\end{nghiepvu}
\goiy{Ngày chơi có doanh thu quảng cáo 0,0031 USD như giá trị mẫu có phải ngày không có doanh thu không?}
\goiy{Đếm có điều kiện bằng \code{COUNT(CASE WHEN ... THEN 1 END)}, so tổng hai cột với 0 mà không làm tròn trước.}
\ghichuthietke{Dễ: 2 bảng, chữ trên đường nối; câu chỉ cần \code{player_daily}. Hai cột doanh thu không có ô trống. Bẫy ở public: làm tròn 2 chữ số trước khi so biến hàng nghìn ngày có doanh thu quảng cáo rất nhỏ thành 0; thiếu cận trên lấy cả tháng 8. hidden-1 bỏ khoảng 10\% người chơi nên số khác public.}
\end{cauhoi}

% ================================================================ N7 · Sửa lỗi · Cơ bản · A2 · after_sales
\begin{cauhoi}[ma=jelly-n7-crash-null, dulieu=jelly, domain=after_sales,
  hinhthuc=debug, cap=basic, dang=A2, kynang={WHERE, so chuỗi}, bay={khoang-trang-thua, nham-nguon-dem},
  tacgia={A}, trangthai=review, bokiemtra={public, hidden-1}, thutu=co]
\tieude{Số lượt chơi gặp sự cố ra 0}
\boicanh[Bạn Vy · Thực tập sinh phân tích]{Em đếm số lượt chơi gặp sự cố trong tháng 7 để gửi nhóm QA thì ra 0, trong khi người chơi báo lỗi khá nhiều. Nhờ anh chị kiểm tra lại giúp em.}
\begin{yeucau}
  \item Sửa truy vấn dưới đây để đếm đúng số lượt chơi có kết quả gặp sự cố trong tháng 7/2026.
  \item Giữ nguyên cột \code{so_luot_su_co}.
\end{yeucau}
\begin{erd}
  \erdbang{level_attempts}{0,0}{attempt_id:PK, event_date*=2026-07-18, level_id, result*='crash', app_version, platform}
\end{erd}
\begin{nghiepvu}{Ước lượng mức độ sự cố trong màn chơi để xếp lịch sửa lỗi.}
  \khainiem{ty-le-su-co}
  \luuy{Cột \code{result} có bốn giá trị viết thường: \code{'complete'}, \code{'fail_out_of_moves'}, \code{'quit'}, \code{'crash'}. Khi so chuỗi, dấu cách ở cuối cũng được tính.}
\end{nghiepvu}
\goiy{So giá trị trong điều kiện lọc với giá trị mẫu của \code{result} trong ERD, từng ký tự.}
\ghichuthietke{Cơ bản: lỗi một chỗ, giá trị mẫu là manh mối. Collation \code{utf8mb4_0900_ai_ci} không phân biệt hoa thường nhưng là NO PAD, nên dấu cách thừa ở \code{'Crash '} làm điều kiện không khớp dòng nào. Bài sai: sửa đúng chuỗi nhưng xoá luôn điều kiện tháng (rớt ở public); lấy \code{SUM(crash_count)} của \code{player_daily} (nguồn đếm khác, lệch xa, rớt ở public). Mã câu giữ theo bản đồ dù bẫy là khoảng trắng, không phải NULL.}
\end{cauhoi}

% ================================================================ N8 · Sửa lỗi · Dễ · A2 · finance
\begin{cauhoi}[ma=jelly-n8-ty-le-sai-chia, dulieu=jelly, domain=finance,
  hinhthuc=debug, cap=easy, dang=A2, kynang={phép tính, ngoặc, ROUND}, bay={thu-tu-phep-tinh, lam-tron-som},
  tacgia={A}, trangthai=review, bokiemtra={public, hidden-1}, thutu=co, saiso=0.0001]
\tieude{ARPDAU tự tính lệch với cột có sẵn}
\boicanh[Bạn Nam · Thực tập sinh phân tích]{Em tính lại ARPDAU từng ngày từ 15/8 tới 21/8 để đối chiếu với bảng chỉ số ngày. Ngày nào số em tính cũng lớn hơn cột có sẵn rất nhiều.}
\begin{yeucau}
  \item Truy vấn dưới đây tính ARPDAU (doanh thu bình quân trên một người chơi hoạt động trong ngày) của từng ngày từ 15/8/2026 tới hết 21/8/2026 nhưng ra số quá lớn. Sửa lại cho đúng công thức ARPDAU.
  \item Giữ nguyên các cột \code{date}, \code{dau}, \code{arpdau_tinh} (USD, làm tròn 4 chữ số), \code{arpdau} và thứ tự sắp.
\end{yeucau}
\begin{erd}
  \erdbang{daily_kpi}{0,0}{date:PK*, dau*=1450, iap_revenue*=120.45, ad_revenue*=20.37, arpdau*=0.0971}
  \erdbang{player_daily}{1,0}{user_id:PK, activity_date:PK, iap_revenue, ad_revenue}
  \erdnoi[kieu=1-n, tren={tổng hợp theo ngày, không có khoá}]{daily_kpi.date}{player_daily.activity_date}
\end{erd}
\begin{nghiepvu}{Kiểm tra bảng chỉ số ngày trước khi đưa ARPDAU vào báo cáo tuần.}
  \khainiem{dau}
  \khainiem{arpdau}
  \luuy{Cột \code{arpdau} có sẵn được làm tròn 4 chữ số; số tự tính cũng phải làm tròn 4 chữ số mới so được với cột này.}
\end{nghiepvu}
\goiy{Trong biểu thức hiện tại, phép chia cho \code{dau} áp dụng cho những số nào?}
\goiy{Đặt \code{iap_revenue + ad_revenue} trong ngoặc trước khi chia cho \code{dau}.}
\ghichuthietke{Dễ: 2 bảng, đường nối ghi rõ là tổng hợp theo ngày. Chấm với \code{saiso=0.0001} vì ARPDAU chỉ cỡ vài phần trăm USD. Starter chỉ chia doanh thu quảng cáo cho DAU rồi cộng cả doanh thu IAP nên số lớn gấp nhiều lần. Bài sai: làm tròn 2 chữ số (rớt ở public với sai số 0,0001); bỏ doanh thu quảng cáo khỏi tử số (rớt ở public). hidden-1 tính lại \code{daily_kpi} nên số khác public.}
\end{cauhoi}

% ================================================================ N9 · Sửa lỗi · Dễ · A3 · customer
\begin{cauhoi}[ma=jelly-n9-nhan-nguoi-tra-tien, dulieu=jelly, domain=customer,
  hinhthuc=debug, cap=easy, dang=A3, kynang={CASE WHEN, so sánh biên}, bay={bien-lon-hon-bang, nham-doanh-thu-quang-cao},
  tacgia={A}, trangthai=review, bokiemtra={public, hidden-1}, thutu=co]
\tieude{Nhóm trả tiền gồm mọi người chơi}
\boicanh[Bạn Long · Thực tập sinh phân tích]{Em chia người chơi hoạt động ngày 15/8 thành nhóm trả tiền và nhóm không trả tiền, nhưng nhóm trả tiền lại gồm mọi người chơi. Nhờ anh chị xem lại điều kiện gán nhãn giúp em.}
\begin{yeucau}
  \item Xét người chơi hoạt động ngày 15/8/2026.
  \item Sửa truy vấn dưới đây để gán nhãn \code{'tra_tien'} cho người chơi có doanh thu IAP (mua hàng trong ứng dụng) trong ngày (\code{iap_revenue}) lớn hơn 0, còn lại gán \code{'khong_tra_tien'}.
  \item Giữ nguyên các cột \code{nhom}, \code{so_nguoi_choi} và thứ tự sắp.
\end{yeucau}
\begin{erd}
  \erdbang{player_daily}{0,0}{user_id:PK, activity_date:PK*=2026-08-15, iap_revenue*=0.00, ad_revenue=0.0125}
  \erdbang{users}{1,0}{user_id:PK, platform, country}
  \erdnoi[kieu=n-1, tren={\code{player_daily.user_id = users.user_id}}]{player_daily.user_id}{users.user_id}
\end{erd}
\begin{nghiepvu}{Đếm người chơi trả tiền trong ngày để theo dõi tỷ lệ người chơi trả tiền.}
  \khainiem{nguoi-choi-tra-tien}
  \luuy{Người chơi chỉ xem quảng cáo tạo ra doanh thu quảng cáo nhưng không phải người chơi trả tiền.}
\end{nghiepvu}
\goiy{Người chơi có \code{iap_revenue} bằng 0 như giá trị mẫu có thoả điều kiện hiện tại không?}
\goiy{Đổi phép so \code{>=} thành \code{>}.}
\ghichuthietke{Dễ: 2 bảng, chữ trên đường nối; giá trị mẫu 0,00 là manh mối. Doanh thu không âm nên starter gán mọi người vào \code{'tra_tien'}. Bài sai cộng cả doanh thu quảng cáo vào điều kiện (nhầm người tạo doanh thu với người trả tiền) rớt ở public. hidden-1 bỏ khoảng 10\% người chơi nên số khác public.}
\end{cauhoi}


% ################################################################ CASE STUDY · CƠ BẢN

% ================================================================ C1 · Case · Cơ bản · finance
\begin{cauhoi}[ma=jelly-c1-quang-cao-co-thuong, dulieu=jelly, domain=finance,
  hinhthuc=case, cap=basic, diem=10, tacgia={A}, trangthai=review, bokiemtra={public, hidden-1}, thutu=co, saiso=0.01]
\tieude{Quảng cáo có tặng thưởng mang lại bao nhiêu?}
\boicanh[Chị Ngọc · Quản lý doanh thu]{Nhóm thiết kế muốn giảm quảng cáo xen kẽ vì người chơi khó chịu. Trước khi quyết định, chị cần biết tháng 7 mỗi loại quảng cáo thu được bao nhiêu.}
\begin{erd}
  \erdbang[hl]{ad_events}{0,0}{ad_id:PK, event_date*=2026-07-21, ad_type*='rewarded', placement, revenue_usd*=0.00612}
\end{erd}
\begin{nghiepvu}{Quyết định có giảm quảng cáo xen kẽ hay không.}
  \khainiem{doanh-thu-quang-cao}
  \luuy{Quảng cáo có tặng thưởng (\code{'rewarded'}) do người chơi tự chọn xem để nhận phần thưởng; quảng cáo xen kẽ (\code{'interstitial'}) tự hiện giữa các màn. Doanh thu trên 1.000 lượt xem dùng để so hai loại có số lượt xem khác nhau.}
\end{nghiepvu}
\begin{buoc}[ma=b1, diem=5, hinhthuc=write, bang={ad_events}]{Doanh thu theo loại quảng cáo}
  \begin{yeucau}
    \item Xét các lượt xem quảng cáo trong tháng 7/2026.
    \item Tính cho từng loại quảng cáo (\code{ad_type}) số lượt xem, doanh thu và doanh thu trên 1.000 lượt xem.
    \item Trả về \code{ad_type}, \code{so_luot_xem}, \code{doanh_thu_usd} (USD, làm tròn 2 chữ số), \code{doanh_thu_1000_luot_usd} (USD, làm tròn 2 chữ số).
    \item Sắp \code{doanh_thu_usd} giảm dần.
  \end{yeucau}
  \goiy{Gom theo \code{ad_type}; doanh thu trên 1.000 lượt bằng 1000 nhân tổng \code{revenue_usd} rồi chia số lượt xem.}
\end{buoc}
\begin{buoc}[ma=b2, diem=5, hinhthuc=ketluan, bang={ad_events}]{Kết luận}
  \begin{luachon}
    \phuongan*{Quảng cáo có tặng thưởng tạo ra gần chín phần mười doanh thu quảng cáo tháng 7 nhờ số lượt xem lớn; doanh thu trên 1.000 lượt của hai loại gần bằng nhau.}
      {Đúng. Bước 1 cho thấy lượt xem có tặng thưởng gấp khoảng tám lần, còn doanh thu trên 1.000 lượt của hai loại chỉ chênh chưa tới 0,1 USD.}
    \phuongan{Mỗi lượt xem quảng cáo có tặng thưởng trả cao gấp nhiều lần quảng cáo xen kẽ, nên đây là loại đáng giá hơn hẳn trên từng lượt xem.}
      {Chưa đúng. Chênh lệch nằm ở tổng doanh thu do số lượt xem; doanh thu trên 1.000 lượt của hai loại gần bằng nhau.}
    \phuongan{Quảng cáo xen kẽ chỉ mang lại phần doanh thu nhỏ, nên bỏ hẳn loại này sẽ không làm doanh thu quảng cáo của game giảm.}
      {Chưa đúng. Bước 1 cho thấy quảng cáo xen kẽ vẫn đóng góp khoảng một phần mười doanh thu quảng cáo của tháng.}
  \end{luachon}
\end{buoc}
\ghichuthietke{Case Cơ bản: 1 bảng, 1 bước SQL, 1 bước kết luận đọc từ kết quả. Doanh thu trên 1.000 lượt là cách so quen thuộc của mạng quảng cáo (eCPM), đề nói bằng lời. Bài sai ở bước 1: thiếu cận trên lấy cả tháng 8; đếm người xem thay lượt xem. Cả hai rớt ở public. Tỷ lệ lượt xem và doanh thu trên 1.000 lượt giữ gần như cũ ở hidden-1. Bẫy kết luận: nhầm tổng doanh thu với doanh thu trên từng lượt; coi phần nhỏ là không đáng kể.}
\end{cauhoi}

% ================================================================ C2 · Case · Cơ bản · sales
\begin{cauhoi}[ma=jelly-c2-gia-goi-ban-chay, dulieu=jelly, domain=sales,
  hinhthuc=case, cap=basic, diem=10, tacgia={A}, trangthai=review, bokiemtra={public, hidden-1}, thutu=co, saiso=0.01]
\tieude{Gói rẻ nhất có bán chạy nhất?}
\boicanh[Anh Huy · Quản lý doanh thu]{Có đề xuất giảm giá thêm cho gói rẻ nhất vì gói này bán chạy. Trước buổi họp, anh cần số giao dịch và doanh thu của từng gói trong hai tháng qua.}
\begin{erd}
  \erdbang[hl]{purchases}{0,0}{transaction_id:PK, event_date*=2026-08-03, product_type*='starter_pack', price_usd*=0.99, is_refund=0}
\end{erd}
\begin{nghiepvu}{Quyết định có giảm giá thêm cho gói rẻ nhất hay không.}
  \khainiem{doanh-thu-iap}
  \luuy{Gói bán chạy là gói có nhiều giao dịch; gói có doanh thu cao là gói có tổng tiền lớn. Hai cách xếp có thể khác nhau khi giá các gói chênh nhau nhiều.}
\end{nghiepvu}
\begin{buoc}[ma=b1, diem=5, hinhthuc=write, bang={purchases}]{Số giao dịch và doanh thu theo gói}
  \begin{yeucau}
    \item Xét các giao dịch từ ngày 1/7/2026 tới hết ngày 29/8/2026.
    \item Tính cho từng gói (\code{product_type}) số giao dịch và doanh thu.
    \item Trả về \code{product_type}, \code{so_giao_dich}, \code{doanh_thu_usd} (USD).
    \item Sắp \code{doanh_thu_usd} giảm dần.
  \end{yeucau}
  \goiy{Gom theo \code{product_type}; số giao dịch là \code{COUNT(*)}, doanh thu là \code{SUM(price_usd)}.}
\end{buoc}
\begin{buoc}[ma=b2, diem=5, hinhthuc=ketluan, bang={purchases}]{Kết luận}
  \begin{luachon}
    \phuongan*{Starter Pack có nhiều giao dịch nhất nhưng doanh thu thấp nhất trong năm gói; Booster Bundle mang về doanh thu cao nhất.}
      {Đúng. Bước 1 cho thấy xếp theo số giao dịch và xếp theo doanh thu cho hai thứ tự khác nhau.}
    \phuongan{Starter Pack bán chạy nhất nên cũng tạo ra doanh thu cao nhất; giảm giá thêm sẽ làm doanh thu của game tăng.}
      {Chưa đúng. Giá Starter Pack thấp nhất nên doanh thu của gói thấp nhất dù nhiều giao dịch; bước 1 cũng không cho biết giảm giá làm số giao dịch tăng bao nhiêu.}
    \phuongan{Season Pass có giá cao nhất nên thu được doanh thu cao nhất; nên dồn ưu đãi cho Season Pass thay cho Starter Pack.}
      {Chưa đúng. Season Pass có ít giao dịch; bước 1 cho thấy Booster Bundle có doanh thu cao hơn.}
  \end{luachon}
\end{buoc}
\ghichuthietke{Case Cơ bản: 1 bảng, 1 bước SQL, 1 bước kết luận. Không cần lọc \code{is_refund} vì trong kỳ chưa có hoàn tiền (khái niệm doanh thu IAP đã nêu). Bài sai ở bước 1: đếm người mua thay giao dịch (rớt ở public); không giới hạn ngày (rớt ở hidden-1, vì biến thể có thêm giao dịch Booster Bundle ngày 2026-09-01). Bẫy kết luận: nhầm ``bán chạy'' (số giao dịch) với doanh thu cao; suy ra tác động của giảm giá khi chưa có số liệu.}
\end{cauhoi}

\end{document}
